\documentclass[12pt]{article}
\usepackage[utf8]{inputenc}
\usepackage[T1]{fontenc}
\usepackage{amsmath, amssymb, geometry, booktabs}
\usepackage{enumitem} % <-- Needed for [label=(\alph*)]
\geometry{margin=1in}
\setlength{\parindent}{0pt}
\renewcommand{\arraystretch}{1.3}


\title{Practice Problem Set 2}
\author{ECON 308 -- International Economic Relations}
\date{Fall 2025}

\begin{document}
\maketitle

\section*{Short Conceptual Questions}

\begin{enumerate}
\item Explain how an increase in a country’s labor endowment affects its production possibilities frontier and pattern of trade under the Heckscher–Ohlin model.

\item Define the \textbf{Rybczynski theorem}. What does it predict about output levels when one factor endowment increases, holding goods prices constant?

\item Suppose capital is relatively scarce in Country A compared to Country B. According to the Stolper–Samuelson theorem, what happens to the real wage and real rental rate in Country A when it opens to trade?

\item Distinguish between \textbf{export-biased growth} and \textbf{import-biased growth}. How does each affect a country’s terms of trade?

\item In the Standard Trade Model, what is meant by an “improvement in the terms of trade”? How does it influence national welfare?

\item Define \textbf{external economies of scale}. How do they differ from \textbf{internal economies of scale} in terms of market structure and trade patterns?

\item Briefly describe how \textbf{path dependence} can lead to persistent differences in industrial location under external economies.
\end{enumerate}

\newpage

\section*{Analytical and Numerical Problems}

\subsection*{Problem 1: Long-Run Factor Adjustment (Heckscher–Ohlin Model) }

Two countries, Home (H) and Foreign (F), produce two goods: textiles (T) and steel (S).  
Textiles are \textbf{labour-intensive}, and steel is \textbf{capital-intensive}.  

\[
K_H = 400, \quad L_H = 800, \quad K_T/L_T = 0.25, \quad K_S/L_S = 2.
\]

\begin{enumerate}[label=(\alph*)]
\item Compute Home’s relative factor abundance.  
\item Which good will Home export under free trade?  
\item Using the Stolper–Samuelson theorem, predict how trade affects real wages and real returns to capital in Home.  
\item Suppose the world price of textiles rises by 10\%. Explain qualitatively how factor prices will adjust.  
\end{enumerate}

\vspace{1em}

\subsection*{Problem 2: The Standard Trade Model and Welfare }

A country produces food (F) and electronics (E) with a concave PPF given by:
\[
F^2 + 2E^2 = 200.
\]
The world relative price is \( P_E/P_F = 2 \).

\begin{enumerate}[label=(\alph*)]
\item Determine the production point that maximizes the value of output.  
\item If consumption preferences imply \( F = E \), find the consumption bundle and trade pattern.  
\item Evaluate qualitatively whether the country gains from trade compared to autarky using indifference curves and the terms of trade concept.  
\item Suppose the terms of trade improve (export price rises). Explain the welfare effect using a diagrammatic argument.  
\end{enumerate}

\vspace{1em}

\subsection*{Problem 3: External Economies and Industrial Location}

Consider two identical countries, A and B, producing semiconductors under \textbf{external economies of scale}.  
Industry output in a country determines unit cost:
\[
C = 20 - 0.05Q,
\]
where \( Q \) is total industry output in that country.  
World demand is:
\[
Q^W = 400 - 5P.
\]

\begin{enumerate}[label=(\alph*)]
\item Derive the relationship between price \( P \) and total world output \( Q^W \).  
\item Suppose initially each country produces half of world output. Compute the equilibrium price.  
\item Show that if one country produces slightly more, its average cost falls—illustrate why this can lead to complete specialization in one location.  
\item Explain how such specialization can result in trade even between identical countries.  
\item Discuss how government policy (e.g., temporary subsidy or protection) could shift production location permanently.  
\end{enumerate}



\end{document}
